\documentclass{article}
\usepackage[utf8]{inputenc}
\usepackage[T1]{fontenc}
\usepackage{lmodern}
\usepackage{hyperref}
\usepackage{geometry}
\geometry{a4paper, margin=1in}

\title{Prompt Engineering Quiz: Week 3 \& 4 (Answer Key)}
\author{Ankara Medipol University}
\date{\today}

\begin{document}

\maketitle

\section*{Answer Key}

\begin{enumerate}

    \item \textbf{c) Role/Audience} \\
    \textit{Explanation: This component sets the persona for the AI and defines who the response is intended for, which controls tone and complexity.}

    \item \textbf{b) Few-shot prompting} \\
    \textit{Explanation: Few-shot prompting provides a small number of examples (shots) to guide the model's output style and format.}

    \item \textbf{b) To ground the model's answers in provided, specific context.} \\
    \textit{Explanation: RAG connects the LLM to an external knowledge source, ensuring its answers are based on that specific data rather than its general training.}

    \item \textbf{b) Zero-shot Chain-of-Thought (CoT)} \\
    \textit{Explanation: This simple phrase encourages the model to "show its work" and output its reasoning process, which often improves accuracy on complex tasks, even without prior examples.}

    \item \textbf{c) They clearly separate instructions from data or context, reducing ambiguity.} \\
    \textit{Explanation: Delimiters help the model understand which part of the prompt is the instruction and which part is the content it should work on.}

    \item \textbf{d) Conflicting directives} \\
    \textit{Explanation: This is a classic anti-pattern where the prompt gives two mutually exclusive instructions, confusing the model.}

    \item \textbf{d) CO-STAR} \\
    \textit{Explanation: CO-STAR is a framework designed to ensure all key components of a good prompt are considered.}

    \item \textbf{b) Frameworks package multiple patterns into a reusable, structured template for a workflow.} \\
    \textit{Explanation: Patterns are the individual building blocks (like role or format), while frameworks are the recipes that combine them for a specific purpose.}

    \item \textbf{b) Meta-Prompting} \\
    \textit{Explanation: Meta-prompting is the practice of using prompts to reason about or generate other prompts.}

    \item \textbf{c) Generate a structured JSON object specifying a function to be called and its arguments.} \\
    \textit{Explanation: The model doesn't execute the tool itself; it signals its intent for your code to run a specific function with specific parameters.}

    \item \textbf{c) Prompt Playbook} \\
    \textit{Explanation: A playbook is the most comprehensive asset, detailing an entire process that includes multiple prompts, human checkpoints, and operational guidelines.}

    \item \textbf{b) Adjust} \\
    \textit{Explanation: The IDEA loop consists of Identify, Draft, Evaluate, and Adjust, representing a full cycle of prompt refinement.}

\end{enumerate}

\end{document}
